\documentclass[12pt,a4paper]{article}
\usepackage[utf8]{inputenc}
\usepackage[french]{babel}
\usepackage[T1]{fontenc}
\usepackage{amsmath, amsfonts, amssymb, graphicx, geometry, xcolor, titlesec, hyperref}

\geometry{margin=2.5cm}
\titleformat{\section}{\Large\bfseries\color{blue!70!black}}{}{0em}{}
\titleformat{\subsection}{\large\bfseries}{}{0em}{}

\title{Rapport de Travaux Pratiques — TP5 \\ \large SVD (Singular Value Decomposition)}
\author{NEUSSI NJIETCHEU PATRICE EUGNE \\ Matricule : 24P820}
\date{Dimanche 5 Avril 2026}

\begin{document}
\maketitle

\section{Introduction}
Ce document présente la réalisation du TP pour la semaine correspondant au module sur \textbf{SVD (Singular Value Decomposition)}.
L'objectif principal était de: Decomposition en valeurs singulieres via Jacobi One-Sided.

\section{Realisation technique}
L'implementation a ete faite en C++17 en suivant les regles du cours (Zero STL, Only Standard Headers).
Le systeme de build \textbf{Jenga} a ete utilise pour la compilation.

\section{Resultats}
Le programme a ete compile et execute avec succes sur Linux.
Les resultats obtenus sont conformes aux attentes theoriques.

\section{Code Source}
Le code source complet est disponible sur le depot GitHub associe : \\
\texttt{https://github.com/PatriceEugene/TP5_SVD_Solvers_Patrice}

\end{document}
